% Bayesian Inference for Parameter Estimation
% Covers: prior distributions, constraints, MC sampling, nested sampling,
%         results analysis, and integration with the sweep pipeline.
% Uses macros from preamble.tex

\section{Bayesian Inference for Parameter Estimation}
\label{sec:inference}

\textbf{Status:} Implemented (Monte Carlo + nested sampling via dynesty/PolyChord) \\
\textbf{Date:} 2026-04-02

Parameter sweeps (\cref{sec:cli-reference}) explore a fixed grid, which scales
exponentially with the number of parameters and does not produce posterior
distributions or evidence estimates.  The \texttt{tidal sample} subcommand adds
a proper Bayesian inference layer: it draws samples from prior distributions,
evaluates each point by running a full PDE simulation, and returns weighted
posterior samples together with Bayesian evidence.

\subsection{Problem Statement}
\label{sec:inference-problem}

Grid sweeps are effective for 1--2 parameters but scale poorly: a
$10\times10\times10$ grid over three coupling constants requires $10^3 =
1{,}000$ simulations, many of which land in unphysical or uninformative
regions.  Bayesian inference solves both problems:

\begin{enumerate}
  \item \textbf{Posterior analysis.}  Monte Carlo sampling assigns probability
    mass according to the prior and likelihood, so the posterior identifies
    which parameter combinations are consistent with observations.

  \item \textbf{Model comparison.}  Nested sampling computes the Bayesian
    evidence $\mathcal{Z} = \int \mathcal{L}(\theta)\,\pi(\theta)\,\mathrm{d}\theta$,
    enabling rigorous model selection via the Bayes factor
    $B_{12} = \mathcal{Z}_1/\mathcal{Z}_2$.
\end{enumerate}

\subsection{Prior Distributions}
\label{sec:inference-priors}

Each parameter requires exactly one prior specified via
\texttt{--prior "NAME=DIST:ARG1:ARG2"}.  Four families are supported:

\begin{table}[htbp]
\centering
\caption{Prior distribution families.}
\label{tab:priors}
\begin{tabular}{@{}llll@{}}
\toprule
Distribution & CLI name & Parameters & Support \\
\midrule
Uniform              & \texttt{uniform}         & $a$, $b$     & $[a,\, b]$ \\
Log-uniform          & \texttt{log\_uniform}    & $a$, $b$ ($a>0$) & $[a,\, b]$ \\
Normal               & \texttt{normal}          & $\mu$, $\sigma$  & $(-\infty,+\infty)$ \\
Arctan-uniform       & \texttt{arctan\_uniform} & (ignored --- warns; use \texttt{0:0}) & $(-19.98,+19.98)$ \\
\bottomrule
\end{tabular}
\end{table}

\paragraph{Arctan-uniform prior}
The \texttt{arctan\_uniform} distribution covers a broad symmetric range
including sign changes --- useful for coupling constants where the sign is
physically significant.  It draws $\theta$ uniformly in
$(-\pi/2 + \epsilon,\, \pi/2 - \epsilon)$ and returns $x = \tan(\theta)$,
giving the fixed support $|x| \le \tan(\pi/2 - \epsilon) \approx 19.98$ with
$\epsilon = 0.05$.  The log-probability is the truncated Cauchy density
$-\log(1+x^2) - \log(\Delta\theta)$.

\textbf{The CLI bounds are ignored for this family} (GH \#425): the support
above is fixed regardless of the \texttt{LO:HI} arguments, and constructing
a prior with bounds that differ from it emits a \texttt{UserWarning}.  Use
the sanctioned sentinel \texttt{0:0} to declare the bounds deliberately
unused.  (Campaign metadata that records \texttt{arctan\_uniform:-89:89}
describes bounds that were never sampled --- this mismatch was one of the
two mechanisms behind the GH \#420 marginal-$D_\mathrm{KL}$ inflation.)

\paragraph{Effective support --- one source of truth}
Because the recorded bounds can describe nothing, no consumer may re-derive
a prior's range from them.  \texttt{tidal.inference.\_prior.effective\_support(%
distribution, low, high)} is the single accessor for \emph{the range that was
actually sampled}: the recorded bounds for uniform and log-uniform, the fixed
$\pm 19.98$ for arctan-uniform, and $(-\infty, +\infty)$ for normal (whose
support is genuinely unbounded --- callers needing a finite plotting range
derive one themselves rather than being handed a silent truncation).  Every
module that answered this question privately answered it wrongly: the
marginal-$D_\mathrm{KL}$ estimator histogrammed a $\pm 20$ posterior over
$(-89, 89)$ (GH \#420) and the \texttt{--full-prior-bounds} corner plot read
the bounds as \emph{degrees} and drew $\pm\tan(89^\circ) \approx \pm 57.3$
panels (GH \#451).  Chains written from v0.49.2 also record
\texttt{effective\_low}/\texttt{effective\_high} in each scalar prior record
(omitted when the support is unbounded, since \texttt{Infinity} is not valid
strict JSON); archived chains, which record only the requested bounds, have
their support reconstructed on read.

Each prior implements three operations required by the sampling pipeline:
\texttt{sample(rng, n)} for Monte Carlo, \texttt{log\_prob(x)} for posterior
weighting, and \texttt{transform(u)} for nested sampling's unit-cube protocol.

\begin{lstlisting}[style=bash]
# Uniform coupling constant
--prior "alpha=uniform:0.01:10"

# Log-uniform mass ratio (positive, varies over orders of magnitude)
--prior "xi=log_uniform:0.001:1.0"

# Normal near known value
--prior "delta=normal:0.0:1.0"

# Broad symmetric range (+-19.98) including sign changes
--prior "g=arctan_uniform:0:0"   # 0:0 = sanctioned "bounds unused" sentinel
\end{lstlisting}

\subsection{Hard Constraints}
\label{sec:inference-constraints}

Physical consistency conditions (ghost-freedom, stability bounds, positivity)
can be enforced as hard priors via \texttt{--constraint "EXPR"}.  Points that
violate a constraint are assigned $\log\pi = -\infty$ immediately, with no
simulation wasted.

Constraint expressions use Python's \texttt{ast} module for safe evaluation ---
there is no call to \texttt{eval()} and no risk of arbitrary code execution.
Supported operations:

\begin{table}[htbp]
\centering
\caption{Operators and functions available in constraint expressions.}
\label{tab:constraint-ops}
\begin{tabular}{@{}ll@{}}
\toprule
Category & Available \\
\midrule
Comparison & \texttt{>}, \texttt{>=}, \texttt{<}, \texttt{<=}, \texttt{==} \\
Arithmetic & \texttt{+}, \texttt{-}, \texttt{*}, \texttt{/}, \texttt{**} \\
Functions  & \texttt{abs()}, \texttt{sqrt()}, \texttt{log()}, \texttt{exp()}, \texttt{sin()}, \texttt{cos()} \\
\bottomrule
\end{tabular}
\end{table}

\begin{lstlisting}[style=bash]
--constraint "xi > 0"
--constraint "alpha > 0"
--constraint "deltam**2 < 2*alpha*xi"   # ghost-freedom condition
\end{lstlisting}

Multiple \texttt{--constraint} flags are combined as logical AND.
Constraints can also be placed in a TOML sweep configuration:

\begin{lstlisting}[style=toml]
[constraints]
expressions = ["xi > 0", "alpha > 0", "deltam**2 < 2*alpha*xi"]
\end{lstlisting}

\subsection{Likelihood Functions}
\label{sec:inference-likelihood}

The likelihood wraps the existing \texttt{\_simulate\_run()} +
\texttt{\_measure\_run()} pipeline from \texttt{tidal/cli/\_sweep.py}.
Simulation failures (solver divergence, NaN) are mapped to
$\log\mathcal{L} = -\infty$.  Three likelihood modes are available:

\begin{table}[htbp]
\centering
\caption{Likelihood modes.}
\label{tab:likelihood}
\begin{tabular}{@{}llll@{}}
\toprule
Mode        & CLI name     & Formula & Use case \\
\midrule
Maximize    & \texttt{maximize}   & $\log\mathcal{L} = $ metric      & Exploratory: find the largest conversion region \\
Gaussian    & \texttt{gaussian}   & $-\tfrac{1}{2}\!\left(\frac{x-\mu}{\sigma}\right)^{\!2}$ & Comparison to data or theoretical prediction \\
Threshold   & \texttt{threshold}  & $0$ if $x \geq x_\mathrm{min}$, else $-\infty$ & Feasibility: is the target observable? \\
\bottomrule
\end{tabular}
\end{table}

CLI syntax: \texttt{--likelihood "METRIC:TYPE:ARGS"}:

\begin{lstlisting}[style=bash]
# Find largest conversion probability (no target)
--likelihood "P_max:maximize"

# Compare to target 0.5 with uncertainty 0.1
--likelihood "P_max:gaussian:0.5:0.1"

# Require minimum conversion (feasibility scan)
--likelihood "P_max:threshold:0.01"
\end{lstlisting}

\subsection{Sampling Methods}
\label{sec:inference-sampling}

\subsubsection{Simple Monte Carlo (\texttt{--method mc})}
\label{sec:inference-mc}

\begin{enumerate}
  \item Draw \texttt{--n-samples} $N$ parameter vectors from the joint prior.
  \item Reject any vector violating a hard constraint (no simulation).
  \item Run a full simulation for each surviving sample and evaluate the log-likelihood.
  \item Return all $N$ samples with their $\log\mathcal{L}$ and $\log\pi$ values.
\end{enumerate}

Monte Carlo is the natural first step: it produces the full joint posterior
$\log p(\theta) = \log\mathcal{L}(\theta) + \log\pi(\theta)$, histograms readily,
and parallelises trivially over samples.  Use it before committing to
nested sampling.

Parallel evaluation uses \texttt{multiprocessing.Pool} with the same
picklable worker pattern as \texttt{tidal sweep --parallel}.

\subsubsection{Nested Sampling (\texttt{--method nested})}
\label{sec:inference-nested}

Nested sampling~\cite{skilling2004nested} is an algorithm that evaluates the
evidence integral $\mathcal{Z} = \int \mathcal{L}(\theta)\,\pi(\theta)\,\mathrm{d}\theta$
by maintaining a set of $N_\mathrm{live}$ live points and iteratively
replacing the lowest-likelihood point with a new draw from the constrained
prior.  The evidence is accumulated as the prior volume shrinks.

\textbf{Two backends} are supported via \texttt{--sampler}:

\begin{table}[htbp]
\centering
\caption{Nested sampling backends.}
\label{tab:nested-backends}
\begin{tabular}{@{}lllll@{}}
\toprule
Backend     & \texttt{--sampler} & Install & Environment & Parallelism \\
\midrule
\textbf{dynesty}~\cite{speagle2020dynesty} & \texttt{dynesty}   & \texttt{pip install tidal[inference]} & Development & \texttt{pool.map} \\
\textbf{PolyChord} & \texttt{polychord} & GitHub + gfortran & HPC (Newton/CSD3) & MPI + OpenMP \\
\bottomrule
\end{tabular}
\end{table}

\textbf{dynesty} is a pure-Python nested sampler installable with a single
\texttt{pip} command and is the recommended backend for development and
testing.  Dynamic nested sampling (\texttt{--dynamic}) allocates additional
live points to regions of high posterior mass, improving parameter estimation
at modest extra cost.

\textbf{PolyChord} (Handley et al.) uses slice sampling with excellent
performance in high-dimensional curved posteriors and full MPI parallelism
on HPC clusters.  It is reserved for production runs on the Newton server
or CSD3 because it requires a Fortran compiler (\texttt{gfortran}).

\textbf{anesthetic}~\cite{handley2019anesthetic} (Handley 2019) handles all
posterior visualization and is backend-agnostic: switching from dynesty to
PolyChord requires only changing \texttt{--sampler} with no change to any
analysis code.

Key stopping parameters:
\begin{itemize}
  \item \texttt{--nlive N} --- number of live points (default 100; increase for multimodal posteriors)
  \item \texttt{--dlogz F} --- stop when estimated remaining evidence $\Delta\log\mathcal{Z} < F$ (default 0.01)
\end{itemize}

\subsection{Results Analysis}
\label{sec:inference-results}

All sampling methods return an \texttt{InferenceResult} dataclass stored in
\texttt{tidal/inference/\_results.py}:

\begin{lstlisting}[style=python]
@dataclass
class InferenceResult:
    samples: NDArray        # (n_samples, n_params)  parameter vectors
    log_likelihood: NDArray # (n_samples,)            log L for each sample
    log_prior: NDArray      # (n_samples,)            log pi for each sample
    param_names: list[str]  # parameter names
    method: str             # "mc" or "nested"
    log_evidence: float | None       # nested sampling only
    log_evidence_err: float | None   # 1-sigma uncertainty on log Z
    weights: NDArray | None          # importance weights
    metadata: dict                   # seed, wall time, n_rejected, ...
\end{lstlisting}

Key analysis methods:

\begin{table}[htbp]
\centering
\caption{\texttt{InferenceResult} analysis methods.}
\label{tab:inference-methods}
\begin{tabular}{@{}ll@{}}
\toprule
Method & Returns \\
\midrule
\texttt{effective\_sample\_size()} & Kish ESS $= (\sum w_i)^2 / \sum w_i^2$ \\
\texttt{best()}                    & MAP estimate $\argmax_\theta\, p(\theta)$ \\
\texttt{posterior\_mean()}         & Weighted posterior mean \\
\texttt{credible\_interval(0.95)}  & Per-parameter 95\% credible intervals \\
\texttt{to\_sweep\_results()}      & Converts to \texttt{SweepResults} for CSV/JSON export \\
\texttt{save(output\_dir)}         & Writes \texttt{results.csv}, \texttt{results.json}, \texttt{inference.json} \\
\bottomrule
\end{tabular}
\end{table}

Output structure:

\begin{lstlisting}
output_dir/
|- results.csv       # self-contained parameter + metric table (via SweepResults)
|- results.json      # same data in JSON
|- inference.json    # evidence log Z, ESS, MAP estimate, credible intervals
|- corner.png        # if --corner was passed
|- trace.png         # if --trace was passed
\end{lstlisting}

\paragraph{Visualization}
\texttt{plot\_corner()} uses anesthetic's \texttt{NestedSamples.plot\_2d()} when
anesthetic is installed, producing publication-quality corner plots with
1-D histograms and 2-D kernel-density marginals.  If anesthetic is absent,
a matplotlib fallback produces a basic scatter corner plot.
\texttt{plot\_trace()} always uses matplotlib.

\subsection{CLI Reference}
\label{sec:inference-cli}

\begin{lstlisting}[style=bash]
tidal sample <json_path> [options]
\end{lstlisting}

\begin{table}[htbp]
\centering
\caption{CLI flags for \texttt{tidal sample}.}
\label{tab:sample-cli}
\begin{tabular}{@{}llll@{}}
\toprule
Flag & Type & Default & Description \\
\midrule
\texttt{--prior SPEC}       & str (repeatable) & (required) & Prior specification \texttt{NAME=DIST:A:B} \\
\texttt{--constraint EXPR}  & str (repeatable) & none       & Hard constraint expression \\
\texttt{--likelihood SPEC}  & str              & (required) & Likelihood \texttt{METRIC:TYPE[:ARGS]} \\
\texttt{--method}           & \texttt{\{mc,nested\}}    & \texttt{mc}    & Sampling method \\
\texttt{--sampler}          & \texttt{\{dynesty,polychord\}} & \texttt{dynesty} & Nested sampling backend \\
\texttt{--n-samples N}      & int              & 1000       & MC sample count \\
\texttt{--nlive N}          & int              & 100        & Nested sampling live points \\
\texttt{--dlogz F}          & float            & 0.01       & Evidence tolerance \\
\texttt{--dynamic}          & flag             & off        & Dynamic nested sampling \\
\texttt{--seed N}           & int              & 42         & RNG seed \\
\texttt{--parallel N}       & int              & 1          & Parallel workers \\
\texttt{--corner}           & flag             & off        & Generate corner plot \\
\texttt{--trace}            & flag             & off        & Generate trace plot \\
\texttt{--output DIR}       & path             & (required) & Output directory \\
\bottomrule
\end{tabular}
\end{table}

All \texttt{tidal simulate} flags (\texttt{--param}, \texttt{--grid-shape},
\texttt{--bounds}, \texttt{--ic}, \texttt{--t-end}, etc.)\ are accepted and
passed through to each simulation run unchanged.

\subsection{Integration with Existing Infrastructure}
\label{sec:inference-integration}

The inference pipeline reuses the sweep framework's simulation machinery
with zero duplication:

\begin{itemize}
  \item \textbf{Simulation:} \texttt{SimulationLikelihood} in
    \texttt{tidal/inference/\_likelihood.py} calls \texttt{\_simulate\_run()} and
    \texttt{\_measure\_run()} directly from \texttt{tidal/cli/\_sweep.py}.
  \item \textbf{Output:} \texttt{InferenceResult.to\_sweep\_results()} converts
    posterior samples to a \texttt{SweepResults} object, enabling CSV/JSON
    export via the same serialization path as \texttt{tidal sweep}.
  \item \textbf{Parallelism:} The same picklable module-level worker pattern as
    \texttt{\_sweep.py}'s \texttt{\_run\_single\_wrapper} is used for \texttt{Pool}-safe
    likelihood evaluation.
\end{itemize}

Install the optional inference dependencies:

\begin{lstlisting}[style=bash]
pip install "tidal[inference]"   # adds dynesty >= 2.1 and anesthetic >= 2.0
\end{lstlisting}

\subsection{Usage Examples}
\label{sec:inference-examples}

\subsubsection{Exploratory Monte Carlo (no target)}

\begin{lstlisting}[style=bash]
tidal sample examples/data/coupled_scattering.json \
  --prior "g0=uniform:0.01:0.9" \
  --prior "mChi2=log_uniform:0.1:10.0" \
  --likelihood "P_max:maximize" \
  --method mc --n-samples 500 --parallel 4 \
  --corner --output results_mc/
\end{lstlisting}

\subsubsection{Gaussian likelihood (comparison to target)}

\begin{lstlisting}[style=bash]
tidal sample examples/data/coupled_scattering.json \
  --prior "g0=uniform:0.01:0.9" \
  --likelihood "P_max:gaussian:0.5:0.05" \
  --method mc --n-samples 1000 \
  --output results_mc_gauss/
\end{lstlisting}

\subsubsection{Nested sampling with ghost-freedom constraints}

\begin{lstlisting}[style=bash]
tidal sample examples/data/torsion_dark_photon.json \
  --prior "alpha=uniform:0.01:10" \
  --prior "xi=log_uniform:0.001:1.0" \
  --prior "deltam=arctan_uniform:0:0" \
  --constraint "xi > 0" \
  --constraint "deltam**2 < 2*alpha*xi" \
  --likelihood "P_max:maximize" \
  --method nested --nlive 200 --dlogz 0.01 \
  --sampler dynesty --parallel 8 --corner \
  --output results_nested/
\end{lstlisting}

Evidence and parameter constraints are written to
\texttt{results\_nested/inference.json}.

\subsubsection{Snapshot count for PolyChord/nested-sampling workloads}
\label{sec:inference-snapshot-guidance}

For PolyChord and dynesty nested-sampling runs where the likelihood is
\texttt{peak\_conversion} (or any measurement whose maximum occurs at
$t = t_\mathrm{end}$ in the perturbative linear regime), pass
\texttt{--snapshots=\$t\_end} to request exactly two snapshots
(initial state + final state). The modal solver's Pad\'{e} path then
applies $\exp(M \cdot t_\mathrm{end})$ in a single step, giving
machine-precision at $t_\mathrm{end}$ and eliminating $N_\mathrm{snaps}-1$
unnecessary IFFT reconstructions per likelihood call. This is also the
most numerically clean configuration: no accumulated matvec roundoff
(see \cref{sec:modal-snapshot-precision}).

\begin{lstlisting}[style=bash]
# 2-snapshot machine-precision inference (t_end=10)
tidal sample spec.json \
  --prior "zeta1=uniform:-0.05:0.05" \
  --likelihood "P_max:maximize" \
  --method nested --sampler polychord --nlive 400 \
  --snapshots 10 \
  --output results_polychord/
\end{lstlisting}

For visualization runs (\texttt{tidal simulate}, \texttt{tidal sweep}),
the default of $t_\mathrm{end}/20$ snapshot interval (21 snapshots) is
appropriate.

\subsection{Parameter Importance via KL Divergence}
\label{sec:inference-importance}

Nested sampling produces weighted posterior samples.  The \emph{joint}
diagnostics are computed by anesthetic; the \emph{per-parameter marginals}
are computed by \textsc{tidal}'s own histogram estimator
(\texttt{tidal.inference.\_importance}, rewritten for GH \#420):

\begin{description}
  \item[$D_\mathrm{KL}$] (anesthetic) Kullback--Leibler divergence from
    prior to posterior (nats).  Measures total information gained.
  \item[$d_G$] (anesthetic) Bayesian model dimensionality: the effective
    number of parameters constrained by the likelihood.  If
    $d_G \ll n_\mathrm{params}$, most parameters are prior-dominated.
  \item[Marginal $D_\mathrm{KL}$] (\textsc{tidal}) Per-parameter KL
    divergence against that parameter's own prior.  Each posterior column
    is first mapped into the space where its prior is uniform ---
    identity (uniform), $\log x$ (log-uniform), $\arctan x$ on the fixed
    $\epsilon$-truncated range (arctan-uniform; the recorded bounds are
    ignored, matching the sampler), the Gaussian CDF (normal) --- and
    histogrammed there against the uniform reference.  Cubed-sphere
    (\texttt{radial\_angular}) couplings have no closed-form 1-D marginal
    CDF and are scored against a fixed-seed empirical prior sample binned
    in $\operatorname{asinh}$ space.
\end{description}

Every computation also emits a \texttt{consistency} self-check block
(recorded in \texttt{importance.json}, \texttt{schema\_version} 2):
superadditivity of marginals against the joint (exact for product priors),
the Kish effective sample size $N_\mathrm{eff}$, the per-parameter noise
floor $(n_\mathrm{bins}-1)/(2N_\mathrm{eff})$ with the
\texttt{floor\_dominated\_params} that sit at or below it, plus
\texttt{fallback\_params} and \texttt{range\_clipped} markers for columns
whose prior record was missing or did not match the samples.  A
\texttt{importance.json} without this block predates v0.48.8 and its
marginals must not be quoted.

\textbf{Usage:}
\begin{verbatim}
# Inline after sampling:
tidal sample spec.json --method nested --analyze ...

# Post-hoc on saved results:
tidal analyze ns_output/ --inference --importance
\end{verbatim}

Interpretation thresholds: $D_\mathrm{KL} > 1$~nat = strong,
$0.1$--$1$~nat = moderate, $< 0.1$~nat = weak --- \emph{subject to the
noise floor}: a marginal at or below its
$(n_\mathrm{bins}-1)/(2N_\mathrm{eff})$ floor is estimator bias, not
constraint, and must not be ranked (GH \#433; the CLI table, the
importance bar chart and the TeX emitters all annotate such values).

\subsection{Finding Extreme Physics}
\label{sec:inference-extremize}

Three likelihood types support directed exploration of amplification:

\begin{description}
  \item[\texttt{maximize}] $\log\mathcal{L} = P_\mathrm{max}$.
    Concentrates samples where amplification $A$ is largest.
  \item[\texttt{minimize}] $\log\mathcal{L} = -P_\mathrm{max}$.
    Finds maximum suppression (smallest $A$).
  \item[\texttt{extremize}] $\log\mathcal{L} = |\log(P / P_\mathrm{GR})|$
    where $P_\mathrm{GR}$ is evaluated per-point from
    \texttt{--baseline-formula}.  Finds both amplification \emph{and}
    suppression simultaneously.
\end{description}

The baseline formula is evaluated with each point's own parameter values,
correctly handling swept $B_0$, $\kappa$, and $t_\mathrm{end}$.

\subsection{High-Dimensional Sampling}
\label{sec:inference-highdim}

The \texttt{--nlive-auto} flag scales live points with dimensionality
following Handley et~al.\ (2015): \texttt{fast} = $\max(25d, 50)$,
\texttt{standard} = $\max(25d, 100)$, \texttt{production} =
$\max(50d, 250)$.  For the 4D dark-photon model this gives 100/100/200;
for 10D non-minimal PGT, 250/250/500.

PolyChord's slice sampling scales better than rejection-based methods
(dynesty/MultiNest) above $\sim$10 dimensions, making it the default
backend.

\subsection{Design Decisions}
\label{sec:inference-design}

\subsubsection{PolyChord as Default Backend}

PolyChord (Handley et~al.\ 2015) is the production nested sampler.
It uses slice sampling (not rejection), handles multimodal distributions
via clustering, and integrates natively with anesthetic.  It requires
\texttt{gfortran}; install via \texttt{bash scripts/install\_polychord.sh}.
dynesty remains available as a pure-Python fallback (\texttt{--sampler dynesty}).

\subsubsection{Why anesthetic for visualization?}

anesthetic (Handley 2019)~\cite{handley2019anesthetic} reads posterior samples
from both dynesty and PolyChord natively.  Using it for all visualization
means the analysis layer is backend-agnostic: switching the sampler from
dynesty to PolyChord requires no change to plotting code.

\subsubsection{Why safe AST evaluation for constraints?}

Using \texttt{eval()} for user-supplied constraint strings would allow arbitrary
Python code execution.  The \texttt{ast} module approach parses expressions into
a syntax tree and evaluates only recognized node types (comparisons, arithmetic,
a fixed set of math functions).  Unknown nodes raise \texttt{ConstraintError}
immediately.  This preserves the flexibility of user-specified expressions
without any security risk.

\subsection{References}
\label{sec:inference-references}

\begin{itemize}
  \item \textbf{Skilling (2004)}~\cite{skilling2004nested}, ``Nested Sampling'',
    AIP Conference Proceedings 735, 395--405. --- Original nested sampling
    algorithm, evidence integral decomposition, convergence criteria.

  \item \textbf{Speagle (2020)}~\cite{speagle2020dynesty}, ``dynesty: a dynamic
    nested sampling package'', MNRAS 493(3), 3132--3158.
    \url{https://github.com/joshspeagle/dynesty} --- The dynesty library
    used as the default development backend; dynamic nested sampling extension.

  \item \textbf{Handley (2019)}~\cite{handley2019anesthetic}, ``anesthetic:
    nested sampling visualisation'', JOSS 4(37), 1414.
    \url{https://github.com/handley-lab/anesthetic} --- Posterior
    visualization library; backend-agnostic, native PolyChord and dynesty support.

  \item \textbf{Barker (2024)}~\cite{barker2024poincare}, ``Poincar\'e gauge
    theory of gravity'', arXiv:2406.12826. --- Identifies Bayesian inference
    of PGT parameters from CMB observations as a natural next step
    (Section~VI.C); the present framework is that future work.
\end{itemize}
